\begin{figure*}[!htb]
\small
\begin{tcolorbox}[
    colback=white, 
    colframe=blue2,
    boxrule=0.5mm,
    arc=0mm,
    outer arc=0mm,
    title=Potential Tools: Operating System - Generation Prompt
]
Task Description:\\
You need to go ahead and generate samples of tool usage with as much diversity as possible, based on the current dimension, the requirements requirements of the task, and the sample references provided.\\

Current dimension:\\
Application Environment Dimension - Special Environment Constraints: Operating System (OS)\\

Task Requirement:\\
- Generate new samples taking into account as much as possible the diversity of user query and available tools.
- Use "// Raw Input //" and "// Raw output //" to label the inputs and outputs of the samples.
- The samples need to consider a realistic scenario as much as possible. Multiple tools need to be available in the list of available tools in the raw input. Please note that in the current dimension, the raw output takes place and occurs in the OS environment in the planning phase, and requires the setting of user query related to the OS level, e.g., kernel updates, user settings, software installation, file or folder handling, boot settings, permissions management for files and directories, etc.\\

At build time, you need to:\\
Start by giving an overall task, then express specific task requirements separately, and make sure that they are as relevant to the os environment as possible. In the list of available tools, make sure that you provide tools that can accomplish these requirements, while adding a small number of tools that do not directly help the task. Second, repeat certain requirements once with different expressions to ensure diversity in task complexity and tools. Finally, complete the corresponding task planning process in the output based on these requirements.\\

Sample Reference:\\
==== Reference Sample Start ====\\
// Raw Input //\\
\colorbox{blue2!50}{\texttt{[sample\_input]}}\\

// Raw output //\\
\colorbox{blue2!50}{\texttt{[sample\_outputl]}}\\
==== Reference Sample End ====\\

Now, please combine the reference sample and the corresponding build steps, try to select tasks from different domains and scenarios, and use different types of tools to make the new sample different from the reference sample in terms of task content and context. We start generating new single samples corresponding to the dimensions.\\
Special Environment Constraints: Operating System (OS) generation sample (one):\\
==== Sample Start ====
\end{tcolorbox}
\caption{Generation prompt for potential tools scenario, OS subtask.}
\label{fig:os-prompt}
\end{figure*}